\section{Глава 2. Проектирование системы}


\subsection{2.1 Пользовательские сценарии}
Проектирование системы опирается на ролевую модель «владелец поездки - участник поездки» и на сквозные сценарии, формирующие целевую пользовательскую ценность. Владелец инициирует поездку и управляет доступом, участники участвуют в обсуждениях, планировании маршрута и распределении расходов.

Ключевые сценарии, учитываемые при проектировании, представлены в таблице.

\begin{table}[H]
\centering
\small
\caption{Ключевые пользовательские сценарии и риски}
\begin{tabular}{|p{3.2cm}|p{3.2cm}|p{4cm}|p{3.4cm}|}
\hline
Сценарий & Компоненты & Результат & Риски \\ 
\hline
Создание поездки и приглашение участника & Клиент Android, API поездок, API приглашений & Формируется общая рабочая область поездки для владельца и участника & Невалидная или устаревшая ссылка приглашения \\ 
\hline
Идея $\rightarrow$ обсуждение $\rightarrow$ маршрут & Подсистема идей, комментарии real-time, модуль маршрута & Предложение участника после обсуждения фиксируется в плане дня & Конкурирующие изменения статуса идеи \\ 
\hline
Добавление расхода и пересчет баланса & Подсистема расходов, участники поездки, клиентский расчет итогов по данным API & Пользователи видят актуальные доли и итоговый баланс & Ошибки входных данных расхода \\ 
\hline
Погода и AI-предложение & Погодный модуль, AI-модуль, подсистема идей & Пользователь получает контекст и сохраняет предложение в идеи & Частичная недоступность внешних сервисов \\ 
\hline
Работа без сети и синхронизация & Локальное хранилище, очередь изменений, API синхронизации & Операции создания, изменения и удаления фиксируются локально и отправляются после восстановления сети & Коллизии при одновременных изменениях \\ 
\hline
\end{tabular}
\end{table}

Для демонстрации связности модулей в главе 2 используется базовый поток «идея -> обсуждение -> маршрут», поскольку именно он отражает переход от коллективного обсуждения к формальному плану поездки.

\textbf{[МЕСТО ДЛЯ РИСУНКА 2.2: ключевой пользовательский поток "идея -> обсуждение -> маршрут" (вставляется вручную)]}
Подпись к вставке: схема потока из 3-5 блоков (\texttt{идея -> обсуждение -> добавление в маршрут}), источник --- текущая версия приложения; рисунок должен подтверждать связность шагов пользовательского сценария.

\subsection{2.2 Архитектура системы}
Система проектируется как клиент-серверное приложение. Клиентская часть реализует пользовательские экраны, локальное хранение и оркестрацию сценариев. Серверная часть централизует бизнес-правила, доступ к данным и интеграцию с внешними сервисами. Такой подход поддерживает единый источник актуального состояния поездки при сохранении отзывчивости интерфейса [12][13][14].

\textbf{[МЕСТО ДЛЯ РИСУНКА 2.1: обзорная архитектура системы (экспорт схемы/скрин архитектурной диаграммы) (вставляется вручную)]}
Подпись к вставке: архитектурная диаграмма \texttt{Android client -> Backend -> PostgreSQL/интеграции}, источник --- проектная схема; рисунок должен подтверждать границы подсистем и направления обмена данными.

\begin{figure}[H]
\centering
\fbox{\parbox{0.94\textwidth}{\textbf{Логика архитектурной схемы приложения:}\\ Android-клиент отправляет запросы в backend API, локально сохраняет изменения при отсутствии сети, после восстановления подключения синхронизируется с сервером. Сервер работает с PostgreSQL и интегрируется с OpenWeather API, YandexGPT API и push-провайдером.}}
\caption{Текстовое представление архитектурной схемы (визуальная схема вставляется вручную)}
\end{figure}

Архитектурная схема фиксирует два уровня согласования данных: оперативный (онлайн-взаимодействие через API и каналы обновлений) и отложенный (офлайн-изменения с последующей синхронизацией). Такое разделение необходимо для мобильного сценария, в котором состояние сети не гарантировано [19].

\subsection{2.3 Модель данных и связи сущностей}
Концептуальная модель данных строится вокруг сущности поездки как корневого объекта. Все функциональные подсистемы связаны с поездкой и наследуют ее контекст (даты, участники, ограничения, валюту).

Основные связи сущностей описываются следующим образом:
1. Пользователь может участвовать в нескольких поездках, а одна поездка содержит нескольких участников.
2. Идея принадлежит конкретной поездке и может быть преобразована в активность маршрута.
3. Комментарии являются элементами обсуждения идеи и связаны с автором.
4. Маршрут представлен набором дней, каждый день содержит упорядоченный список активностей.
5. Расход относится к поездке и включает распределение долей между выбранными участниками.
6. Погодные данные привязываются к выбранному городу и датам поездки.
7. Предложение AI создается в контексте поездки и может быть сохранено как идея.

Для серверного хранения применяется реляционная модель данных на основе PostgreSQL. В текущей реализации структура таблиц инициализируется схемой проекта и поддерживается в кодовой базе. На стороне клиента сохраняется локальное состояние, необходимое для восстановления пользовательского процесса при временной потере сети [10][11][19].

\subsection{2.4 Подсистемы и функциональное назначение}
Для проектирования и последующей реализации выделены функциональные подсистемы, каждая из которых решает отдельную задачу, но работает в рамках общего сценария поездки.

\begin{table}[H]
\centering
\small
\caption{Подсистемы приложения и их назначение}
\begin{tabular}{|p{2.7cm}|p{3.3cm}|p{3.3cm}|p{3.3cm}|}
\hline
Подсистема & Функции & Пользовательская ценность & Ограничения \\ 
\hline
Авторизация и профиль & Вход, управление профилем, выход, удаление аккаунта & Контролируемый доступ к персональным данным & Зависимость от доступности внешней аутентификации \\ 
\hline
Поездки и участники & Создание/редактирование поездки, управление участниками, приглашения & Быстрое формирование совместного пространства & Валидация дат и проверка корректности приглашений \\ 
\hline
Идеи и обсуждения & Создание идей, комментарии, модерация статуса & Прозрачная фиксация коллективных решений & Ролевые ограничения операций владельца/участника \\ 
\hline
Маршрут по дням & Добавление, перенос и упорядочивание активностей & Формирование исполнимого плана поездки & Ограничение действиями в диапазоне дат поездки \\ 
\hline
Расходы и баланс & Добавление расходов, сплиты, контроль статусов оплаты & Прозрачность финансовой ответственности & Валидность сумм, состава участников и валюты \\ 
\hline
Погода и AI-подсказки & Прогноз по городу поездки, генерация и сохранение предложений & Поддержка принятия решений по активности & Зависимость от внешних API и качества входного контекста \\ 
\hline
Настройки уведомлений & Управление категориями уведомлений & Снижение риска пропуска важных изменений & Ограничения платформенной доставки push \\ 
\hline
Офлайн и синхронизация & Локальная запись, отложенная отправка операций создания/изменения/удаления & Непрерывность работы при нестабильной сети & Конфликтные изменения при поздней синхронизации \\ 
\hline
\end{tabular}
\end{table}

Представленная декомпозиция соответствует требованиям ТЗ и задает основу для модульной реализации, где пользователь наблюдает целостный процесс, а не набор изолированных экранов [7].

\subsection{2.5 Выбор методов и средств реализации}
Выбор методов и средств реализации выполнен по явным критериям: (1) поддержка мобильного сценария и офлайн-работы, (2) предсказуемость сопровождения, (3) совместимость со стеком Android/Backend, (4) стоимость доработок при расширении продукта.

\subsubsection{2.5.1 Серверная часть}
Серверная часть ориентирована на централизованную реализацию доменных правил и внешних интеграций. Выбран стек Kotlin + Ktor [12][13] по следующим причинам:
1. единый язык с клиентом снижает когнитивные и эксплуатационные издержки команды;
2. Ktor предоставляет REST и WebSocket в одном фреймворке, что покрывает сценарии API и обсуждений;
3. стек допускает модульную структуру маршрутов и тестирование без избыточного фреймворк-каркаса.

В качестве альтернатив рассматривались Java/Spring и Node.js/Express; для данного проекта выбран Kotlin/Ktor как более согласованный со стеком Android и учебным контуром сопровождения.

\subsubsection{2.5.2 Хранилище данных}
Для транзакционного хранения доменных сущностей используется PostgreSQL [10][11]. Выбор обусловлен:
1. реляционной моделью с явными связями между поездками, участниками, идеями, маршрутом и расходами;
2. ACID-семантикой для операций, влияющих на совместное состояние группы;
3. воспроизводимым развертыванием через инициализационную схему проекта.

В качестве альтернатив рассматривались встраиваемые и документные хранилища; на текущем этапе они не обеспечивают такой же предсказуемости транзакционных операций для выбранной модели данных.

\subsubsection{2.5.3 Клиентская часть}
Мобильный клиент реализуется для платформы Android на Kotlin с Jetpack Compose, Hilt и Retrofit [14][15][16]. Выбор обоснован:
1. Compose --- декларативное управление состоянием сложных экранов и быстрые итерации UI;
2. Hilt --- управляемая сборка зависимостей для модульных экранов;
3. Retrofit/OkHttp --- типизированные клиентские контракты и стандартный стек сетевого взаимодействия.

Альтернативный XML-based UI для текущего состава экранов не выбран из-за большей стоимости поддержки состояния и изменений интерфейса.

\subsubsection{2.5.4 Интеграционные компоненты}
Для расширения пользовательского контура в проекте используются внешние сервисы прогноза погоды и AI-генерации. В подсистеме офлайн-работы применяется механизм фоновой обработки задач синхронизации. Такой выбор позволяет сохранить работоспособность ключевых сценариев даже при частичной недоступности внешних источников [17][18][19].

\subsection{2.6 Нефункциональные решения}
Нефункциональные решения сформулированы с позиции эксплуатационной устойчивости и качества пользовательского опыта.

1. \textbf{Согласованность данных.} Серверный контур обеспечивает единые правила валидации и хранения, а также разрешение коллизий при синхронизации, чтобы все участники поездки видели согласованное состояние после операций создания, изменения и удаления [12][13].
2. \textbf{Надежность пользовательских сценариев.} Критические действия (идеи, маршрут, расходы) проектируются с учетом обработки ошибок и восстановления после сетевых сбоев [7][19].
3. \textbf{Безопасность доступа.} Доступ к данным поездки ограничен аутентификацией и участием пользователя в поездке; ролевые операции владельца и участника разграничены в сценариях модерации и управления доступом [7][12].
4. \textbf{Интеграционная устойчивость.} При недоступности отдельных внешних сервисов базовые функции совместного планирования остаются доступными, а вспомогательные данные обновляются после восстановления интеграций [17][18].
5. \textbf{Эксплуатационная применимость.} Архитектура допускает развертывание серверной части в облачном контуре и публикацию мобильного клиента в стандартном канале распространения Android-приложений [20][21].

Для проверки проектных решений, связанных с мобильной эксплуатацией, в текст включается отдельная иллюстрация офлайн-режима и восстановления согласованного состояния.

\textbf{[МЕСТО ДЛЯ РИСУНКА 2.3: офлайн-режим и синхронизация изменений (вставляется вручную)]}
Подпись к вставке: схема/скрин сценария \texttt{действие офлайн -> запись в очередь -> восстановление сети -> синхронизация}, источник --- экран приложения или схема из проекта; рисунок должен подтверждать непрерывность сценария при потере сети.

\subsection{2.7 Выводы по главе}
Во второй главе выполнено проектирование системы на уровне пользовательских сценариев, архитектуры, модели данных и функциональной декомпозиции. Показано, что выбранные проектные решения ориентированы на поддержку полного цикла совместного планирования поездки в едином мобильном приложении.

Сформированы проектные основания для реализации ключевых модулей: авторизации, управления поездкой и участниками, идей и обсуждений, маршрута, расходов, погодного и AI-контекста, уведомлений, офлайн-работы и синхронизации. Обоснован выбор методов и средств реализации, достаточный для перехода к этапу программной реализации, представленному в следующей главе.
